Para planificar una estación de acabado se aproxima el número de piezas pendientes de completar por una cantidad continua. Al inicio hay \(12\) piezas pendientes y llegan \(6\) piezas por minuto. La capacidad nominal es de \(10\) piezas por minuto cuando hay trabajo suficiente. No se conocen los tiempos individuales ni cómo cambia el caudal real cuando quedan pocas piezas. Los números son sintéticos.

Construya un modelo para decidir si la estación llega a quedarse sin trabajo durante los primeros diez minutos. Defina qué incluye su estado, complete las hipótesis que faltan y resuelva el modelo elegido. Explique qué hará su modelo si el estado llega a cero. Valore si la capacidad nominal basta para asegurar su conclusión y describa una observación de planta que resolvería la incertidumbre principal.
